% \subsection{Text-to-CAD Generation}
% \paragraph{Text-to-CAD Generation.}
% Learning-based CAD generation has evolved from structured design modeling to language-conditioned generation with pretrained models. Early work modeled CAD construction histories, sketch constraints, and assembly relations using datasets such as DeepCAD~\citep{wu2021deepcad}, SketchGraphs~\citep{seff2020sketchgraphs}, and Fusion 360 Gallery~\citep{willis2021fusion,willis2022joinable}. Task-specific transformers have subsequently been developed to map text or engineering drawings to CAD command sequences~\citep{khan2024text2cad,qin2025drawing2cad}. Pretrained LLMs and VLMs have also been adapted to generate code-like or structured CAD representations~\citep{li2025cad,govindarajan2025cadmium,yuan2024openecad}. Programmatic representations such as CadQuery support directly executable generation~\citep{xie2025text}, while tool-augmented, self-refinement, and evolutionary approaches improve generated CAD programs through iterative feedback~\citep{mallis2025cad,li2026seek,preintner2025evocad}. Our work studies Text-to-CAD generation under joint geometric and physical requirements, including global dimensional bounds and FEM-verified physical constraints.

% \paragraph{Physics-Aware CAD Generation and Engineering Agents.}
% Recent work has started to evaluate generated CAD models against engineering requirements rather than geometric similarity alone. \citep{berger2026physics} integrate physics-based verification into a multi-agent CAD design and refinement loop. \citep{son2026self} study CAD generation from single-part and multi-part engineering briefs, using structured blueprints, visual inspection, and finite-element feedback to iteratively refine generated STEP models. COSMO-Agent~\citep{deng2026cosmo} uses tool-augmented reinforcement learning to train an LLM for closed-loop CAD--CAE optimization, where the agent updates geometric parameters and material choices for parameterized components according to simulation feedback. Related LLM-based systems automate workflows in computational fluid dynamics~\citep{chen2024metaopenfoam,xu2025cfdagent,yue2025foam}, finite-element analysis~\citep{hou2025autofea,qi2025feagpt}, and multiphysics simulation~\citep{zhang2025mooseagent}. 
% Complementary to iterative CAD--CAE refinement, PhyCAD focuses on improving the engineering feasibility of designs at generation time. By learning from CAD--CAE verification outcomes across design tasks, the model is explicitly optimized to jointly satisfy geometric specifications and physical constraints when generating executable CadQuery programs. Engineering feedback thus improves the generator's ability to produce compliant designs, without requiring a simulation-guided refinement loop at inference time.

% \paragraph{Reinforcement Learning with Verifiable Feedback.}
% RL with verifiable feedback provides a practical way to optimize language models when correctness can be evaluated automatically. GRPO~\citep{shao2024deepseekmath} enables policy optimization using group-relative rewards, while InternBootcamp~\citep{li2025internbootcamp} provides diverse tasks with programmatic verifiers. In CAD, CAD-Coder~\citep{guan2026cad} combines supervised fine-tuning with GRPO using geometric and format rewards for Text-to-CAD generation, while cadrille~\citep{kolodiazhnyi2025cadrille} uses geometric feedback for multimodal CAD reconstruction. ReCAD~\citep{li2026recad} combines parameterized-code guidance with hierarchical primitive learning, organizing training over increasingly complex geometric structures. These methods primarily optimize geometric fidelity and output validity rather than FEM-verified physical constraint satisfaction. PhyCAD instead incorporates CAD--CAE feedback into Text-to-CAD policy training through a Hierarchical Balanced Advantage. Unlike a curriculum over geometric primitives, our hierarchy organizes candidates by verification status, combining between-level progress with objective-wise peer ranking within each level.


\paragraph{Text-to-CAD Generation.}
Early learning-based CAD methods modeled construction histories, sketch constraints, and assembly relations using datasets such as DeepCAD~\citep{wu2021deepcad}, SketchGraphs~\citep{seff2020sketchgraphs}, and Fusion 360 Gallery~\citep{willis2021fusion,willis2022joinable}. Subsequent approaches map text or engineering drawings to CAD sequences using task-specific transformers~\citep{khan2024text2cad,qin2025drawing2cad}, or adapt pretrained LLMs and VLMs to structured CAD generation~\citep{li2025cad,govindarajan2025cadmium,yuan2024openecad}. CadQuery supports executable program generation~\citep{xie2025text}, while tool augmentation, self-refinement, and evolutionary search enable iterative refinement~\citep{mallis2025cad,li2026seek,preintner2025evocad}. PhyCAD studies generation under joint geometric and physical requirements, including global dimensional bounds and FEM-verified physical constraints.

\paragraph{Physics-Aware CAD Generation and Engineering Agents.}
Recently, \citep{berger2026physics} integrates physics-based verification into multi-agent CAD refinement, while \citep{son2026self} combines structured blueprints, visual inspection, and finite-element feedback to refine single-part and multi-part STEP models. COSMO-Agent~\citep{deng2026cosmo} trains an LLM through tool-augmented RL to optimize parameterized components through CAE-guided updates to geometry and material choices. Related agents automate computational fluid dynamics~\citep{chen2024metaopenfoam,xu2025cfdagent,yue2025foam}, finite-element analysis~\citep{hou2025autofea,qi2025feagpt}, and multiphysics simulation~\citep{zhang2025mooseagent}. Complementary to these engineering workflows, PhyCAD addresses a learning challenge in physics-aware CAD generation: prioritizing joint geometric and physical feasibility while retaining informative feedback on partial progress. We introduce a Hierarchical Balanced Advantage that combines progress across verification levels with objective-wise peer ranking within each level. This organizes heterogeneous CAD--CAE outcomes into a training signal that preserves cross-level priority while guiding constraint satisfaction and improvements within the feasible design space.

\paragraph{Reinforcement Learning with Verifiable Feedback.}
GRPO~\citep{shao2024deepseekmath} enables policy optimization using group-relative rewards, while InternBootcamp~\citep{li2025internbootcamp} provides diverse tasks with programmatic verifiers. In CAD, CAD-Coder~\citep{guan2026cad} combines supervised fine-tuning and GRPO using geometric and format rewards, while cadrille~\citep{kolodiazhnyi2025cadrille} uses geometric feedback for multimodal reconstruction. ReCAD~\citep{li2026recad} combines parameterized-code guidance with hierarchical primitive learning, following a curriculum over increasingly complex geometric structures. These methods primarily target geometric fidelity and output validity rather than FEM-verified physical constraints. PhyCAD incorporates CAD--CAE feedback into policy training through a Hierarchical Balanced Advantage. Rather than scheduling a geometric curriculum, it groups candidates by verification status and combines between-level progress with objective-wise peer ranking within each level.